\scanpage{34}
% L1–5 | top, title and definition
\uline{Catégorie des algèbres bivariantes} (sur un anneau $k$)

C'est la catégorie dont les objets sont les $k$-algèbres\inserted{\uline{\unsure{muni|munies|munie} d'une forme} $k_A:A\unsure{-|\to} k$}, et les morphismes de $A$ dans $B$ les couples $f=(f_*,f^*)$, où $f^*:B-A$ est un hom d'algèbres, et où $f_*:A-B$ est un hom $B$-linéaire\inserted{[\unsure{à|:} g. et \unsure{à|:} dr]}, i.e. un hom additif satisfaisant les ``form. de projection''

% L6–8 | upper, projection formulas and annotation
\[
\bigl[f_*(x f^*(y))=f_*(x)y\bigr],\quad f_*(\unsure{f|f^*}(y)x)=y f_*(x). \tag*{$(\ast)$}
\]
\margin{signe $(-1)^{\delta j}$ \unsure{raisonnable} ds le cas gradué, si la deux. formule est valable, $y$ de degré $j$, $f_*$ de degré $\delta$.}
% slanted
\margin{Il faut choisir entre l'une et l'autre formule !}
% slanted
et de plus :
\[
k_B f_*=k_A \tag*{$(\ast\ast)$}
\]

% L9–13 | upper middle
On a une variante graduée, \deleted{et} On aura à travailler dans diverses sous-catégories pleines, ou même non pleines.

Dans le cas gradué anticommutatif (où $f^*$ est donc un hom d'Alg. graduées, et \deleted{\unsure{de}} $f_*$ préserve à translation près \inserted{\uline{de degré pair}} les graduations), les deux formules $(\ast)$ sont équivalentes.

% L13 margin | left, slanted block
\margin{\uline{Remarque} \unsure{Pour|Bien} que $(f_*,f^*)$ soit un isomorphisme, il faut \unsure{et} que $f^*$ \deleted{soit un isom et} $f_*$ \unsure{bijectifs|bijectif}, i.e. $f^*$ un isom et $f_*(1)$ inversible ; les ``bons'' isom. sont ceux où $f_*(1)=1$, ou encore $f_*$ et $f^*$ sont inverses l'un de l'autre.}
% slanted

% L14–19 | middle
Soit $A$ une algèbre munie d'une forme
\[
k_A:A-k
\]
\deleted{que} \unsure{donc|ds|dans} la forme bilinéaire $\unsure{b|B}$
\[
(x,y)\longmapsto k(xy)
\]
sur $A\times A$ \deleted{est séparée, de} défini\unsure{t|\deleted{t}} un \inserted{(= deux !)} hom. \deleted{injectif}
\[
\varphi_A:A\longrightarrow\check A\qquad(\check A:\mathrm{Hom}_{k\text{-}\mathrm{mod}}(A,k)).
\]
\inserted{$\langle y,\varphi_A(x)\rangle=k(yx)$}
\margin{\deleted{\unsure{l'injection} \illegible{2}}}
% slanted

% L20–23 | lower middle
Supposons que $B$ soit de \deleted{\unsure{a}} muni de $k_B:B-k$, \inserted{tel que $\varphi_B$ \uline{soit injectif},} et soit $f^*:B\longrightarrow A$ tel que l'hom transposé $\,{}^{t}f^*:\check A-\check B$ applique $\mathrm{Im}\,\varphi_A$ dans $\mathrm{Im}\,\varphi_B$ [ce qui n'est pas une restriction si $\varphi_B:B\xrightarrow{\sim}\check B$ !)

Alors l'application induite, \inserted{notée} $f_*$, est donc caractérisée par
\[
k(y f_*(x))=k(f^*(y)x) \tag*{$(\ast\ast\ast)$}
\]

% L24–27 | foot, several additions along the lower edge
et satisfait de ce fait la $2^{\text{ème}}$ relation $(\ast)$, [donc la première dans le cas commutatif, ou gradué anticommutatif, quand $f_*$ est hom à translation \inserted{paire} près].
\inserted{De plus, il satisfait $(\ast\ast)$ (faire $y=1$ dans $(\ast\ast\ast)$).}
\deleted{\unsure{Supposons donc} $A,B$}
Réciproquement, $A$ étant \unsure{fixé}, \deleted{\unsure{étant}} tout hom $(f^*,f_*)$ de $(A,k_A)$ en $(B,k_B)$, $f_*$ satisfait $(\ast\ast\ast)$ ce qui prouve \unsure{que $k$ étant donné|$k$ étant donné}, $f_*$ est déterminé par $f^*$.
% archive sheet number: 33

